\chapter{Étape 2 — Séparer les os}
\label{chap:segmentation}

\begin{chapeau}
La segmentation décide, pour chaque voxel, à quel os il appartient. Le résultat est un tableau d'étiquettes : 0 pour
« pas d'os », puis un numéro par os. Les noms (fémur, tibia\ldots) viennent à l'étape 3.
\end{chapeau}

\section{La difficulté : les os se touchent}

Un simple seuil ne suffit pas. Au seuil de l'os (226~HU), tout le membre forme un seul bloc : dans chaque articulation,
quelques voxels de volume partiel relient les deux os (\cref{fig:segmentation-principe}, image 2). Monter le seuil
sépare les os, mais efface aussi l'os spongieux et les corticales fines. L'idée de la méthode est de se servir des
deux seuils : un seuil \emph{haut} pour trouver une graine par os, puis le seuil \emph{bas} pour reconstruire chaque os
en entier.

\fig{segmentation_principe.png}{La segmentation sur une vraie coupe du genou du lévrier : (1) l'image du scanner ;
(2) au seuil de 226~HU, les os se touchent ; (3) avec un seuil plus haut, il reste un noyau par os ; (4) chaque os
grandit jusqu'à l'articulation. La figure est en 2D ; le logiciel travaille en 3D.}{fig:segmentation-principe}

\section{Préparer le volume}

\subsection{Recadrer}

Le volume complet contient la table du scanner, l'air et les tissus mous. Le logiciel cherche d'abord la boîte qui
contient les os denses ($\geq 500$~HU, morceaux d'au moins 30~mm³), en travaillant sur un volume deux fois plus petit
dans chaque direction, puis l'agrandit de 6~mm. Toute la suite se fait dans cette boîte : la mémoire nécessaire est
divisée par cinq environ.

\subsection{Lisser}

Le bruit (environ 49~HU) crée des voxels isolés au-dessus du seuil. On applique un flou gaussien de
$\sigma = 0{,}35$~mm, c'est-à-dire on remplace chaque valeur par une moyenne pondérée de ses voisines :
\begin{equation}
  I_\sigma(\vb{x}) = \int I(\vb{y})\, \frac{1}{(2\pi\sigma^2)^{3/2}} \exp\!\left(-\frac{\norm{\vb{x}-\vb{y}}^2}{2\sigma^2}\right) \dd\vb{y}.
\end{equation}
Le $\sigma$ est donné en millimètres : le flou est le même dans les trois directions, même si les voxels ne sont pas
cubiques. Ce lissage divise le bruit environ par deux sans effacer les interlignes articulaires (1~mm et plus).

\subsection{Le masque osseux}

Le masque osseux est l'ensemble $M = \{\vb{x} : I_\sigma(\vb{x}) \geq 226~\mathrm{HU}\}$. Ce seuil est celui que
propose le logiciel Mimics pour l'os au scanner ; le document de cadrage du projet donnait 200 à 300~HU.

\section{Une graine par os : le seuil qui monte}
\label{sec:graines}

On monte le seuil par paliers de 100~HU, de 500 à 1500~HU. À chaque palier, on regarde les \emph{composantes
connexes} : les groupes de voxels au-dessus du seuil qui se touchent (par une face, une arête ou un coin : 26
voisins). Quand un bloc se coupe en morceaux, c'est en général parce qu'une articulation, moins dense que l'os, vient
de céder.

Un bloc $B$ est coupé au palier $T$ quand, parmi les composantes de $B \cap \{I_\sigma \geq T\}$, au moins deux
\begin{itemize}
  \item sont assez grosses : au moins 300~mm³ ;
  \item et \textbf{persistent} : elles gardent au moins 50~mm³ au palier $T + 200$~HU.
\end{itemize}
Chaque morceau retenu est alors examiné à son tour aux paliers suivants. C'est un arbre : le membre se coupe en
segments, qui se recoupent en os. Les feuilles de l'arbre sont les graines.

\begin{pourquoi}
\textbf{Pourquoi la persistance ?} Un vrai os a une corticale dense, qui reste visible quand le seuil monte encore. Un
fragment d'os spongieux, lui, disparaît vite : sans ce critère, l'épiphyse spongieuse d'un os pourrait devenir une
graine à part.
\end{pourquoi}

\begin{exemple}
Sur le lévrier, le genou se sépare dès 600~HU, la hanche vers 900~HU (la tête du fémur reste collée au bassin plus
longtemps), et les os du tarse vers 1200 à 1400~HU.
\end{exemple}

\section{La croissance de région compétitive}
\label{sec:croissance}

Les graines grandissent ensuite toutes en même temps dans le masque osseux $M$. On utilise pour cela la
\textbf{ligne de partage des eaux à partir de marqueurs} \cite{meyer1994}. Imaginons l'image comme un relief :
\begin{equation}
  R(\vb{x}) = \clamp\bigl(3000 - I_\sigma(\vb{x}),\ 0,\ 4500\bigr).
\end{equation}
Les zones denses sont des vallées, les interlignes (peu denses) des crêtes. Chaque graine est une source d'eau. L'eau
monte : à chaque instant, on inonde le voxel le plus bas du relief (le plus dense) parmi ceux qui touchent une zone
déjà inondée, et il prend la couleur de cette zone. Deux os qui grandissent se rencontrent donc là où la densité est
la plus faible entre eux : dans l'interligne.

\noindent\begin{minipage}{\linewidth}
\begin{lstlisting}[language=Python, title={Principe (pseudo-code)}]
file = voisins des graines, classés par relief croissant
tant que la file n'est pas vide :
    x = le voxel de plus petit relief de la file
    etiquette[x] = etiquette du voisin déjà inondé qui l'a atteint
    ajouter à la file les voisins de x qui sont dans le masque et pas encore inondés
\end{lstlisting}
\end{minipage}

Le « fond » (tout ce qui est hors du masque) est aussi une source : il empêche les os de déborder du masque. Le calcul
est fait par la fonction \code{MorphologicalWatershedFromMarkers} de SimpleITK \cite{lowekamp2013}.

\section{Recoller les morceaux d'un même os}
\label{sec:cstar}

À haut seuil, une épiphyse se détache parfois de la diaphyse, parce que la corticale de la métaphyse est fine : un
même os donne alors deux graines. Mais ces deux morceaux se touchent sur \emph{toute la section} de l'os, alors que
deux os voisins ne se touchent que sur une petite surface d'articulation. On mesure cela par un nombre sans
dimension :
\begin{equation}
  \tcbhighmath{c^* = \frac{A_{\text{contact}}}{V_{\min}^{2/3}},}
  \label{eq:cstar}
\end{equation}
où $A_{\text{contact}}$ est la surface de contact entre les deux morceaux (le nombre de faces de voxels qu'ils
partagent, multiplié par l'aire d'une face) et $V_{\min}$ le volume du plus petit des deux.

\begin{pourquoi}
\textbf{Pourquoi la puissance 2/3 ?} Pour un objet de taille $L$, le volume varie comme $L^3$ et une section comme
$L^2$, donc comme $V^{2/3}$. Le rapport $c^*$ ne dépend donc pas de la taille du chien : une coupure à travers un os
donne toujours $c^*$ de l'ordre de quelques unités, quelle que soit la taille de l'os.
\end{pourquoi}

\fig{critere_contact.pdf}{Le critère $c^*$ mesuré sur le lévrier : les vraies articulations restent sous 2,1 ; un
même os coupé en deux donne environ 5. Le seuil de recollage est fixé à 3.}{fig:critere-contact}

Le recollage est glouton : on fusionne d'abord la paire de plus grand $c^*$, puis on met à jour la table des contacts
sans tout recalculer, grâce à la relation exacte
$A_{\text{contact}}(A \cup B, C) = A_{\text{contact}}(A, C) + A_{\text{contact}}(B, C)$. On recommence tant qu'une
paire dépasse le seuil. Enfin, les segments de moins de 50~mm³ sont rattachés au voisin avec lequel ils ont le plus
grand contact (ou supprimés s'ils sont isolés), et les étiquettes sont renumérotées par volume décroissant.

\section{Résultat}

\fig[0.82]{segmentation_resultat.png}{Les os segmentés automatiquement dans le scanner hanche--genou du
lévrier.}{fig:segmentation-resultat}

\begin{exemple}
Série hanche--genou (728 coupes) : lissage 7~s, graines 10,5~s, croissance 15~s, recollage 5~s. Le fémur fait
100,3~cm³. Pour les séries genou--cheville et pied, la segmentation prend 7 à 8~s (après recadrage).
\end{exemple}

\begin{danslecode}
\fichier{code_source/etape2_segmentation.py} : \code{recadrage()}, \code{lisser()},
\code{graines_automatiques()} (l'arbre des seuils, fonction interne \code{couper()}), \code{croissance()} (ligne de
partage des eaux), \code{contacts()} (surfaces de contact), \code{fusionner_surdecoupages()} (critère $c^*$),
\code{nettoyer()}, et \code{segmenter()} qui enchaîne le tout. Deux outils servent à l'étape 3 :
\code{recroissance()} (nouvelle croissance dans une zone, avec des graines imposées) et \code{raffiner()} (graines plus
fines dans un seul segment). Réglages : \code{SEUIL_OS_HU}, \code{SIGMA_LISSAGE_MM}, \code{SEUIL_GRAINE_*},
\code{PERSISTANCE_HU}, \code{RAPPORT_CONTACT_FUSION}\ldots{} (\cref{ann:parametres}).
\end{danslecode}
